\chapter{Scenario Projection} \label{Chap6}
\section{CMIP6-Derived Climate Scenarios}

% Future weather conditions were obtained from the transient climate dataset developed by \cite{semenov2025transient}. The dataset uses the LARS-WG 8.0 stochastic weather 
% generator to downscale projections from five CMIP6 Global Climate Models for representative locations across Great Britain, including North Wyke. It provides daily 
% minimum and maximum air temperature, precipitation and solar radiation from 2020 to 2090. This study uses the North Wyke simulations under two emissions pathways: the 
% intermediate SSP2-4.5 pathway and the very-high-emissions SSP5-8.5 pathway. Although the source dataset extends to 2090, the present analysis is restricted to 2050.

% The source dataset contains 100 stochastic realisations for each GCM--SSP combination. To retain variation between climate models while keeping the scenario experiment 
% computationally tractable, ten realisations were selected from each of the five GCMs. This produced 50 climate trajectories under each pathway and 100 trajectories 
% overall. Each trajectory contains four directly simulated variables: daily minimum temperature, maximum temperature, precipitation and solar radiation.

% These data differ substantially from the historical NWFP dataset used to develop the forecasting model. The NWFP data comprise tower- and catchment-level observations 
% derived from multiple monitoring systems, whereas the simulated climate data provide site-level daily weather trajectories shared across the three towers. They do not 
% contain FCH\textsubscript{4}, soil conditions, hydrological measurements, livestock records or field-management events. Consequently, the scenario model was restricted to 
% predictors that could either be obtained directly from the climate simulations or reconstructed consistently from the historical record.

% \begin{table}[htbp]
% \centering
% \small
% \renewcommand{\arraystretch}{1.15}
% \begin{tabularx}{\textwidth}{lXX}
% \toprule
% Property & Historical NWFP data & Simulated climate data \\
% \midrule
% Data source
% & Instrumental observations
% & LARS-WG simulations derived from five CMIP6 GCMs \\

% Temporal resolution
% & Hourly observations aggregated to daily values
% & Daily simulations \\

% Spatial representation
% & Tower- and catchment-specific
% & Site-level North Wyke climate \\

% Study period
% & Historical record ending in 2024
% & Future trajectories extending to 2050 \\

% Climate pathways
% & Not applicable
% & SSP2-4.5 and SSP5-8.5 \\

% Available variables
% & Meteorological, soil, hydrological, livestock, management and flux variables
% & Minimum temperature, maximum temperature, precipitation and solar radiation \\

% FCH\textsubscript{4}
% & Observed or gap-filled target
% & Not available; projected by the model \\
% \bottomrule
% \end{tabularx}
% \caption{Comparison between the historical NWFP observations and the simulated climate data used for scenario projection.}
% \label{tab:historical_scenario_comparison}
% \end{table}

% Before constructing the scenario predictor frames, the four simulated climate variables were standardised against the local NWFP climate record. Because the climate 
% trajectories represent the North Wyke site rather than individual tower locations, Tower~4 was used as the common observational reference owing to its comparatively 
% complete meteorological record. Separate correction factors were estimated for each GCM using its near-term SSP2-4.5 simulations from 2020--2030.

% Additive offsets were applied to minimum and maximum temperature, while multiplicative factors were applied to precipitation and solar radiation. The latter preserves 
% non-negative values and is more appropriate for ratio-scale variables. The resulting correction for each GCM was applied unchanged to both SSP2-4.5 and SSP5-8.5. This 
% aligned the simulated variables with the local climate baseline without removing the relative differences between GCMs or emissions pathways.

% The corrected variables were subsequently mapped to the forecasting data structure. Mean air temperature was calculated from the corrected daily minimum and maximum 
% temperatures, precipitation was retained as a daily total, and solar radiation was converted from MJ\,m\textsuperscript{-2}\,day\textsuperscript{-1} to mean 
% W\,m\textsuperscript{-2}. Calendar and historically derived farm-management predictors were then added to form the scenario-compatible input dataset.

Future weather conditions were obtained from the transient climate dataset developed by \cite{semenov2025transient}. It provides daily simulations of minimum and maximum air temperature, precipitation
and solar radiation for North Wyke from 2020 to 2090, generated using five CMIP6 Global Climate Models. This study selected ten realisations per GCM under SSP2-4.5 and SSP5-8.5, producing 100 climate
trajectories, and restricted the projection horizon to 2050.

Unlike the historical NWFP data, these simulations provide site-level daily weather rather than tower- and catchment-specific observations. They contain neither FCH\textsubscript{4} nor the soil,
hydrological, livestock and management measurements available historically. The scenario model was therefore restricted to predictors that could be obtained from the climate simulations or reconstructed
consistently from the NWFP record. Variables unavailable from CMIP6 were reconstructed from historical NWFP
  seasonality. Livestock densities were represented using tower-specific
  day-of-year seasonal mean, while grazing and fertiliser variables were derived
  from historical presence and event patterns. These profiles were repeated
  across future years and modified according to each scenario.

The four simulated climate variables were bias-corrected against the local NWFP record using Tower~4 as the common reference. Per-GCM correction factors were estimated from the 2020--2030 SSP2-4.5
simulations: additive offsets were applied to minimum and maximum temperature, while multiplicative factors were applied to precipitation and solar radiation. The same corrections were applied under both
pathways to preserve their relative differences. Mean temperature was subsequently derived from the corrected minimum and maximum values, and solar radiation was converted from
MJ\,m\textsuperscript{-2}\,day\textsuperscript{-1} to W\,m\textsuperscript{-2}.

Table~\ref{tab:climate_standardisation} reports the resulting correction factors. Raw simulations were cooler than the local record by 1.81--2.71$^\circ$C for daily minimum temperature and
2.60--3.56$^\circ$C for daily maximum temperature. Solar radiation required a modest upward adjustment of 3.4--8.3\%, whereas precipitation was reduced to 22.5--24.0\% of its raw value, indicating that
the uncorrected simulations were approximately 4.2--4.4 times wetter than the NWFP reference.

\begin{table}[htbp]
\centering
\small
\renewcommand{\arraystretch}{1.1}
\begin{tabular}{lrrrr}
\toprule
GCM & $\Delta T_{\min}$ ($^\circ$C) & $\Delta T_{\max}$ ($^\circ$C) &
SWIN ratio & Precipitation ratio \\
\midrule
ACCESS-ESM1-5   & 1.81 & 2.60 & 1.034 & 0.225 \\
CNRM-CM6-1      & 2.26 & 3.09 & 1.083 & 0.229 \\
HadGEM3-GC31-LL & 1.97 & 2.79 & 1.044 & 0.236 \\
MPI-ESM1-2-LR   & 2.71 & 3.56 & 1.060 & 0.226 \\
MRI-ESM2-0      & 2.29 & 3.03 & 1.043 & 0.240 \\
\bottomrule
\end{tabular}
\caption{Per-GCM correction factors used to standardise the simulated climate variables against the Tower~4 NWFP reference. Temperature offsets were added to raw simulated values; solar-radiation and
precipitation ratios were multiplied by raw simulated values.}
\label{tab:climate_standardisation}
\end{table}

\section{Results}

\subsection{Less Informed Model Creation}

This chapter's scenario model uses this dissertation's best forecasting model but restricted to
the 13 features that CMIP6 climate scenario can actually supply (climate drivers, calendar
encodings, livestock density, management practices) rather than the full
52-feature configuration used elsewhere in this dissertation (Chapter~\ref{Chap5}). This
restriction carries an accuracy cost that has been quantified:

\begin{table}[h]
\centering
\small
\renewcommand{\arraystretch}{1.1}
\begin{tabular}{lcc}
\toprule
Configuration & Feature count & MASE (seasonal mean) \\
\midrule
Full-feature model (Chapter~\ref{Chap5}) & 52 & 0.6908 \\
\textbf{Restricted model} & \textbf{13} & \textbf{0.7383} \\
\bottomrule
\end{tabular}
\caption{Accuracy cost of restricting the feature set to what CMIP6
climate scenario can supply. MASE scaled against the day-of-year seasonal mean baseline.}
\label{tab:sp_less_informed}
\end{table}

A restricted model is created using only the features available from the simulated climate
dataset. This model has a 6.9\% increase in MASE relative to the full-feature model (MASE
0.7383 vs.\ 0.6908), but remains sufficiently accurate for a projection task. This restricted
model is used for the inference task in the sections that follow.

% Restricting the feature set to only what a genuine climate scenario can supply carries a small,
% quantifiable accuracy cost: a real-anchor backtest found the direct-regression, per-tower,
% trend-augmented configuration recovers most of that cost (MASE 0.738, vs.\ 0.7405 for this
% dissertation's full-feature TabICLv2 configuration and 0.760 for a superseded one-shot engine) -- a
% $\sim$0.3\% cost rather than the superseded engine's 2.5\%. The trend covariate's own training range
% is $\sim$2.7$\times$ shorter than the horizon it is queried to (2050); a direct sanity check found it
% saturates to a bounded, plausible value rather than diverging when queried decades past its training
% maximum, unlike SARIMAX's known explosive extrapolation behaviour under a comparable stress test.

\subsection{Livestock Density Sensitivity}

Livestock density was varied along a four-level ladder: the historical baseline; half of the
baseline density; a policy-informed reference level of 2.5~LSU\,ha\textsuperscript{-1}; and each
tower's own historical maximum. The reference level is drawn from UK Countryside Stewardship
guidance for non-Severely Disadvantaged Area land \cite{defra2022stewardship}.

Reported scenario responses are mean percentage changes in annual predicted
FCH\textsubscript{4} relative to the matched baseline. They are averaged across the five GCMs and
ten selected realisations per GCM.

\begin{table}[h]
\centering
\small
\renewcommand{\arraystretch}{1.1}
\begin{tabular}{llrrr}
\toprule
SSP & Livestock level & T2 & T4 & T9 \\
\midrule
SSP2-4.5 & Half (0.5$\times$), all species & $-27.5\%$ & $-14.9\%$ & $-26.9\%$ \\
SSP2-4.5 & Government reference (2.5~LSU/ha) & $+4.6\%$ & $+0.9\%$ & $-2.6\%$ \\
SSP2-4.5 & Own historical maximum, all species & $+101.2\%$ & $+138.3\%$ & $+184.3\%$ \\
\midrule
SSP5-8.5 & Half (0.5$\times$), all species & $-27.2\%$ & $-15.0\%$ & $-26.9\%$ \\
SSP5-8.5 & Government reference (2.5~LSU/ha) & $+4.6\%$ & $+0.9\%$ & $-2.6\%$ \\
SSP5-8.5 & Own historical maximum, all species & $+99.7\%$ & $+138.1\%$ & $+181.9\%$ \\
\bottomrule
\end{tabular}
\caption{Mean percentage change in annual predicted FCH\textsubscript{4} relative to the matched
baseline livestock scenario, averaged across five GCMs and ten realisations per GCM.}
\label{tab:livestock_scenario_results}
\end{table}

  Livestock density produced the largest modelled response in the scenario grid.
  Halving stocking density reduced predicted annual FCH\textsubscript{4} by
  15--28\%, while raising it to each tower's historical maximum increased
  predicted flux by 100--184\%. The policy-informed reference level of
  2.5~LSU\,ha\textsuperscript{-1} was close to the historical baseline at T2 and
  T4, but below the baseline at T9, producing only small changes. Responses at
  each livestock level were similar under SSP2-4.5 and SSP5-8.5, indicating that
  the modelled sensitivity to livestock density was comparatively stable across
  the two climate pathways. Illustrative monthly trajectories are provided in
  Appendix Figures~\ref{fig:app_sp_livestock_trajectory_ssp245}
  and~\ref{fig:app_sp_livestock_trajectory_ssp585}.

% Tower 9's negative response at the literature-ceiling rung reflects that its real historical stocking
% density already exceeds that ceiling -- consistent with NWFP's identity as an intensively-managed
% research platform. At the "own historical maximum" rung, Tower 2's response ($+98.7\%$) is decisively
% weaker than Towers 4 or 9's ($+137.8\%$, $+182.9\%$), consistent with Tower 2's established
% livestock-blindness (Chapters~\ref{Chap4},~\ref{Chap5}) -- independently confirmed here for a fourth
% time, under a different model and climate driver source. In absolute annual generation terms, Tower
% 9's own-historical-max scenario reaches 260.7 kg~ha\textsuperscript{-1}~yr\textsuperscript{-1},
% nearly three times its baseline (27.7/72.4/92.2 kg~ha\textsuperscript{-1}~yr\textsuperscript{-1} at
% T2/T4/T9 respectively).

\FloatBarrier

\subsection{Grazing and Fertiliser Sensitivity}

\begin{table}[h]
\centering
\small
\renewcommand{\arraystretch}{1.1}
\begin{tabular}{llrrr}
\toprule
SSP & Intervention & T2 & T4 & T9 \\
\midrule
SSP2-4.5 & Extend grazing season by 4 weeks & $-3.8\%$ & $+19.9\%$ & $+20.0\%$ \\
SSP2-4.5 & Fertiliser frequency $+50\%$ & $-8.4\%$ & $\approx 0.0\%$ & $\approx 0.0\%$ \\
SSP2-4.5 & Fertiliser rate $+50\%$ & $-1.7\%$ & $-2.0\%$ & $\approx 0.0\%$ \\
\midrule
SSP5-8.5 & Extend grazing season by 4 weeks & $-3.7\%$ & $+19.9\%$ & $+19.9\%$ \\
SSP5-8.5 & Fertiliser frequency $+50\%$ & $-8.2\%$ & $-0.1\%$ & $\approx 0.0\%$ \\
SSP5-8.5 & Fertiliser rate $+50\%$ & $-1.7\%$ & $-2.0\%$ & $\approx 0.0\%$ \\
\bottomrule
\end{tabular}
\caption{Mean percentage change in annual predicted FCH\textsubscript{4} relative to the matched
historical scenario, averaged across five GCMs and ten realisations per GCM.}
\label{tab:management_scenario_results}
\end{table}

  Extending the grazing season by four weeks increased predicted annual
  FCH\textsubscript{4} by approximately 20\% at T4 and T9, whereas T2 produced a
  small negative response of approximately 4\%. Fertiliser interventions produced
  smaller and less consistent responses: increasing application frequency reduced
  predicted flux by approximately 8\% at T2 but had little effect at T4 and T9,
  while increasing the application rate changed predictions by no more than 2\%.
  The similarity between SSP2-4.5 and SSP5-8.5 indicates that these modelled
  management sensitivities were stable across the two climate pathways.
  Illustrative trajectories are provided in Appendix
  Figures~\ref{fig:app_sp_grazing_trajectory_ssp245}--%
  \ref{fig:app_sp_fertilizer_trajectory_ssp585}.

% Extending the grazing season by four weeks increases FCH\textsubscript{4} substantially at T4/T9
% ($\approx+20\%$), while T2's response is small and negative ($-3.7\%$ to $-3.8\%$), consistent
% with its muted livestock sensitivity. Fertiliser interventions are comparatively minor, changing
% predicted FCH\textsubscript{4} by at most 8.4\% (T2) and no more than 2\% at T4/T9. As with the
% livestock ladder, the SSP pathway makes almost no difference here -- management practice, not
% climate divergence, drives the response. Trajectories under both pathways are in Appendix
% Figures~\ref{fig:app_sp_grazing_trajectory_ssp245}--\ref{fig:app_sp_fertilizer_trajectory_ssp585}.

% Extending the grazing season produces a real, monotonic increase at Towers 4 and 9; Tower 2's small
% and sign-inconsistent response is consistent with its established livestock-blindness rather than a
% genuinely different underlying finding. Fertiliser perturbations produce a small, sign-inconsistent
% effect at every tower, consistent with this project's standing conclusion that fertiliser timing
% carries little independent signal once livestock density is already present. The SSP2-4.5 vs.\
% SSP5-8.5 pathway divergence is small relative to every management perturbation above throughout most
% of the horizon.

% \section{Limitations} \label{sp:lim}

% Three limitations apply specifically to this chapter's projections, beyond the general absence of
% ground truth. First, TabICLv2's extrapolation behaviour under genuine distribution shift (livestock
% density, climate drivers) has not been formally stress-tested to the standard applied to the
% production forecasting ensemble; the livestock-ladder sweep is a behavioural probe, not a controlled
% extrapolation test with any ground truth to check against. Second, no diagnostic available to this
% dissertation can detect a smooth-but-confidently-wrong extrapolation: historical PICP measures
% calibration only at real, historically-anchored points and cannot be computed for a genuine future
% scenario. Third, the CMIP6-vs-NWFP historical climate mismatch is corrected for but not eliminated as
% a source of residual uncertainty in every downstream scenario result.
